\chapter[Ehrengedächtniß des P. Konrad Georg Engelder]{Ehrengedächtniß des am 2. April 1903 entschlafenen treuverdienten P. Konrad Georg Engelder}
\quellseite{3}

P. Konrad Georg Engelder wurde geboren am 1. März 1837 auf dem St. Leonhardshof bei Beerfelden, Kreis Erbach, Großherzogthum Hessen-Darmstadt. Sein Vater, der ein Tuchmachermeister war, hieß Johann Wilhelm Engelder und seine Mutter Elisabeth, eine geborene Haupt. Er selbst lernte draußen die Weberei. Im Jahre 1854, im Alter von siebzehn Jahren, wanderte er nach America aus und kam zuerst nach Philadelphia, wo er sich auch etliche Jahre aufhielt. Hier arbeitete er eine Zeitlang am Schuhmacherhandwerk.

Engelders deutschländischer Pastor war Rationalist. Die Schule in seinem Heimathsort, welche er besuchte, ließ auch viel zu wünschen übrig; doch lernte er als Knabe Luthers Kleinen Katechismus. Lernbegierig, wie er war, forschte er zu Hause fleißig in der biblischen Geschichte und las begierig alle Welt- und Kirchengeschichtsbücher, deren er habhaft werden konnte. Ueberhaupt trieb er Geschichte mit großer Vorliebe und eignete sich im Laufe der Zeit ein seltenes kirchen- und weltgeschichtliches Wissen an, aus welchem Schatz reichlich mitzutheilen er jederzeit bereit war. Er kam nicht mit der Absicht nach America, Theologie zu studiren, sondern Verdienstes wegen. Anfänglich besuchte er in Philadelphia die Kirche des P. Demme, wo er sich aber nicht heimisch fühlte. Endlich wurde er mit P. Hoyer bekannt, welcher ihn veranlaßte, in unser Predigerseminar zu Fort Wayne einzutreten. Im Herbst des Jahres 1857 bezog er das Seminar. Hier betrieb er nun in rechter Gottesfurcht, treulich und mit großem Fleiß sein Studium unter Anleitung unserer theuren, jetzt seligen Lehrer.

Schon zu Ostern 1860 bestand er mit dreien seiner Studiengenossen sein Abgangsexamen, und man sandte ihn, eher, als ihm lieb war, ins Amt. Er wäre lieber noch länger im Seminar geblieben, um noch mehr zu lernen und noch gründlicher auf die Führung des heiligen Predigtamtes vorbereitet zu werden, fügte sich aber dem Urtheil seiner Professoren. Er folgte einem an ihn ergangenen Ruf zweier lutherischen Gemeinden im Staate New York, zu Olean und Wellsville, in deren Mitte er am 5. Juni 1860 vom seligen P. Dulitz ordinirt und eingeführt wurde. Diesen beiden Gemeinden stand er neun Jahre lang in dem Herrn treulich vor. Im Jahre 1869 erhielt er einen Beruf an die neugegründete lutherische Zions-Gemeinde in Pittsburg, Pa., den er auch mit Einwilligung und nach friedlicher Entlassung seiner bisherigen Gemeinden annahm. Hier wirkte er durch Gottes Gnade in großem Segen bis in das Jahr 1877. Von dort führte ihn der Herr der Kirche an die lutherische Gemeinde in Punxsutawney, Pa., die mit ihrer Filiale in Brady zur Ohio-Synode gehörte, und er wurde somit durch Annahme dieses Berufs Glied der genannten Synode, welche zu jener Zeit mit unserer Synode in Lehr-, Bekenntniß- und Glaubenseinigkeit stand und, wie sie, der Synodalconferenz gliedlich angehörte. Als aber im Jahre 1881 die Ohio-Synode unsere schrift- und bekenntnißgemäße Lehre von der Gnadenwahl verwarf und als calvinistisch bezeichnete, sich auch öffentlich von der Synodalconferenz lossagte, trat unser lieber Engelder mit manchen andern seiner Synodal- und Amtsbrüder gewissenhalber aus der
\quellseite{4}
Ohio-Synode aus, nachdem sie zuvor ernstlich, aber leider vergeblich vor ihr wider ihre falsche Lehrstellung gezeugt und die reine Lehre treulich bekannt hatten. Diese Ausgetretenen bildeten darauf mit ihren Gemeinden die Concordia-Synode, die sich jedoch schon nach vierjährigem Bestand wieder auflöste und in der Missouri-Synode aufging. So wurde Engelder wieder Glied unserer Synode; auch seine beiden Gemeinden schlossen sich ihr an.

Bis zum Jahre 1891 bediente er diese Gemeinden, da berief ihn die evangelisch-lutherische St. Lucas-Gemeinde zu Hannastown, Butler Co., Pa. Er erkannte diesen Beruf als einen göttlichen und folgte ihm. Zwölf Jahre lang hat er dann dieser seiner lieben Gemeinde und unserer Concordia-Waisenheimath, deren Seelsorger er war, wie auch zwei kleinen Gemeinden, die er im Laufe der Jahre noch hinzu bekam, gewissenhaft und segensreich bedient, als ein durch Gottes Gnade frommer, weiser und eifriger Knecht Christi, bis ihn der Herr nach seinem wunderbaren, aber allezeit gnädigen Rath plötzlich mitten aus seiner gesegneten Arbeit zu sich in die ewige, selige Ruhe rief. Noch am Tage vor seinem Tode hielt er Confirmandenunterricht, arbeitete seine Confirmationsrede für den nächsten Sonntag, Palmarum, aus und verrichtete noch andere Werke seines Amtes. Den ganzen Tag war er anscheinend wohl, und nachdem er den Abendgottesdienst mit den Seinen gehalten hatte, begab er sich zur Ruhe. Kaum hatte er eine Stunde geschlafen, da erwachte er, von Athemnoth gequält, eine Folge von Asthma, womit er bereits seit längerer Zeit behaftet war, stand auf und setzte sich an den warmen Ofen. Seine treue Lebensgefährtin stand ihm dienend zur Seite und versuchte, wie sie es schon oft gethan hatte, durch allerlei Mittel sein Leiden zu lindern. Allein diesmal wollte alles nichts helfen; Schmerzen und Angst mehrten sich. Da blickte er sein liebes Weib noch einmal an, die ihren Arm um ihn geschlungen und sein Haupt an ihre Brust gelegt hatte, und mit dem brünstigen Seufzer: „Lieber Gott, hol mich heim!“ entschlief er selig in ihren Armen, auf seinem Stuhle sitzend. Seine Seele, aus aller Angst und Noth gerissen, ist nun daheim bei dem Herrn. Seine irdische Wallfahrt währte 66 Jahre, 1 Monat und 1 Tag. Im Weinberge des Herrn durfte er als sein berufener Diener 43 Jahre arbeiten und hat darin redlich des Tages Last und Hitze getragen, hat aber dabei auch viel Gnade und Segen vom Herrn empfangen.

Im Herbst des Jahres 1860 war er in den heiligen Ehestand getreten mit seiner ihm überlebenden Gattin Katharina, geborenen Graf. Gott segnete diese Ehe mit 14 Kindern, von welchen drei dem Vater in den Himmel vorangegangen sind. Er hinterläßt seinen einzigen hochbetagten Bruder, seine betrübte Wittwe, sieben Söhne, von denen die beiden ältesten im heiligen Predigtamt stehen, vier Töchter, drei Schwiegersöhne, vier Schwiegertöchter und 14 Enkelkinder.

Seine irdische Hülle wurde am 4. April unter großer Betheiligung der Glieder seiner Gemeinden und der Insassen des Waisenhauses als ein edles Samenkörnlein zur Erde bestattet. Alle seine Kinder von nah und fern hatten sich zu seinem Begräbniß eingefunden. Im Hause wie auch am Grabe amtierte P. Lobte, und in der Kirche hielt der Unterzeichnete die Leichenpredigt über Hos. 6, 1.

Wir haben alle an unserm lieben Engelder viel verloren: seine Gattin ihre beste irdische Stütze, seine Kinder den treuen, liebevollen Vater, seine Gemeinde den ihr in dem Herrn wohl vorstehenden, in der reinen Lehre fest gegründeten, gewissenhaften und überaus fleißigen Seelsorger und wir, seine Amtsbrüder, sonderlich wir Glieder der Pittsburg-Conferenz, eines unserer ältesten und erfahrensten Glieder, der in den letzten Jahren unser Vorsteher war, unsere ganze Synode aber einen frommen und brünstigen Fürbitter, der stets innigen Antheil nahm an ihrem und der ganzen Kirche Wohl und Wehe. Ja, es war der Entschlafene ein von Herzen gläubiges, aufrichtiges, demüthiges Gotteskind, ein treuer, fleißiger und geschickter Prediger und Seelsorger, der nicht das Seine, sondern nur Gottes Ehre und das wahre Heil der ihm anvertrauten Seelen suchte; er war uns allen ein Vorbild in seinem Glauben und Wandel. Das glaubte er freilich von sich selber nicht, sondern hielt sich vielmehr für den Geringsten unter uns; er wußte von sich nichts anderes, als daß er ein armer sündhafter Mensch sei, dem aber ohne all sein Verdienst große Barmherzigkeit von Gott widerfahren sei, weil er durch Gottes Gnade Christum als seinen Heiland erkannt habe und von Herzen an ihn glaube und also gewißlich in ihm habe die Erlösung durch sein Blut, nämlich die Vergebung seiner Sünden und das ewige Leben. Und das hat er nun erlangt; weil Christus sein Leben war, ist auch Sterben sein Gewinn gewesen. Sein letzter Wunsch: „Lieber Gott, hol mich heim!“ ist nun erfüllt. Er ist daheim bei seinem Heiland, der ihn von allem Uebel erlöst und ihm ausgeholfen hat zu seinem himmlischen Reich. Dessen wollen wir uns trösten in unserer Trauer über das Hinscheiden unsers geliebten Vaters und Bruders und Gottes Gnade preisen, die sich so herrlich an ihm erzeiget hat. Durch dieselbe Gnade wolle aber der treue Gott auch uns helfen, daß, wenn unser Stündlein vorhanden ist, unser Ende sei wie dieses Gerechten Ende, um Jesu willen. H. St.
